\chapter{Heart Disease Risk Assessment}
\section*{What Is This Test?} Heart-disease risk assessment combines health history, examination, and selected measurements to estimate future cardiovascular risk.\section*{Alternative Names} Cardiovascular risk assessment; ASCVD risk assessment; heart-health check.\section*{Principle} Age, blood pressure, cholesterol, diabetes, smoking, family history, and other factors are integrated in validated risk tools.\section*{Why This Test Is Ordered} It guides prevention decisions in people without known cardiovascular disease and identifies modifiable risk factors.\section*{Primary vs Secondary Test} History, blood pressure, and lipid testing are foundational; ECG, imaging, and biomarkers are not routine screening for everyone.\section*{How This Test Is Performed} The clinician measures blood pressure, reviews health behaviors and medicines, and orders lipid and glucose testing when indicated.\section*{What Preparation Is Needed} Fasting depends on the lipid panel. Bring family history and a complete medicine list.\section*{Understanding Results} A calculated risk is an estimate, not a diagnosis or guarantee; risk tools vary by population and guideline.\section*{Follow-up Tests} Repeat lipids, A1C, ECG, coronary calcium CT, or specialist assessment may be considered based on risk and symptoms.\section*{References} \begin{enumerate}\item MedlinePlus. Heart Disease Risk Assessment.\item American College of Cardiology and American Heart Association. Cardiovascular prevention guideline.\end{enumerate}\clearpage\section*{Understanding Results and Follow-up}
The laboratory report should identify the specimen, method, units, reference interval or decision limit, and whether the result is positive, negative, abnormal, or inconclusive. Reference intervals vary with age, sex, pregnancy, specimen type, assay, and laboratory; a result outside a range is not automatically a diagnosis, and a result within a range does not always exclude disease. Discuss unexpected results with the ordering clinician, who can decide whether repeat or confirmatory testing, treatment, or referral is appropriate. Do not change medicines or delay urgent care based on an isolated result.
\section*{Regulatory Status and Authorization Date}This chapter describes a test category, not a specific commercial product. FDA, CDSCO, EU IVDR, and MHRA approval or registration dates therefore cannot be assigned without the assay manufacturer, product name, instrument, and intended use. For laboratory-developed tests, an individual agency approval date may not exist. See the Preface for the interpretation of regulatory status.
\clearpage
